% ------------------------------------------------------------------------
% file `revisions-exercise-16-exercise-body.tex'
%   in folder `xsim/'
%
%     exercise of type `exercise' with id `16'
%
% generated by the `exercise' environment of the
%   `xsim' package v0.20c (2021/02/03)
% from source `revisions' on 2021/12/14 on line 548
% ------------------------------------------------------------------------
  Résous les problèmes suivants. Laisse ta démarche visible.
  \begin{enumerate}
    \item Dans la cour de récréation, 20 élèves doivent se partager 302 billes. Ali, un élève du groupe, dit : \og Partagez-vous équitablement le maximum de billes, je prendrai celles qui restent ! \fg. Détermine le nombre de billes qu'Ali recevra. Écris tous tes calculs.
    \item À l'entrainement, trois cyclistes font des tours d'un étang. Jean effectue un tour en 9 minutes, Eva en 10 minutes et Philippe en 15 minutes. Ils ont commencé leur entrainement au même endroit et en même temps à 14h15. Détermine l'heure à laquelle ils vont se retrouver à nouveau ensemble à leur point de départ. Écris tous tes calculs.
    \item Ronald et Tim ont fait leur lessive aujourd'hui. Or Ronald fait sa lessive tous les 6 jours et Tim tous les 9 jours. Combien se passera-t-il de jours avant que Ronald et Tim ne refassent leur lessive le même jour ? Laisse ta démarche visible.
    \item Pour une activité, un enseignant répartit 132 filles et 84 garçons en formant le plus grand nombre de groupes mixtes. Tous les élèves participent. Chaque élève n'appartient qu'à un seul groupe. Le nombre de filles est le même dans chaque groupe et le nombre de garçons est le même dans chaque groupe.
          \begin{enumerate}[label=\arabic*)]
            \item Détermine le plus grand nombre de groupes mixtes formés.
            \item Détermine le nombre de filles dans chaque groupe.
            \item Détermine le nombre de garçons dans chaque groupe.
          \end{enumerate}
    \item Sabrina a une chambre rectangulaire qui mesure 5,94 m sur 3,63 m. Elle
          souhaite couvrir le sol de plaques de bois carrées. Réponds aux deux
          questions suivantes en laissant ta démarche visible.
          \begin{enumerate}[label=\arabic*)]
            \item Quelle est la taille maximale possible pour les plaques de bois (en
                  nombre entier de cm) ?
            \item Combien de plaques va-t-elle utiliser ?
          \end{enumerate}
    \item Dans une division euclidienne, le diviseur est 5 et le quotient est 12. Quels sont les dividendes possibles ? Justifie et écris tous tes calculs.
    \item Prouve que la somme de deux nombres pairs consécutifs est un nombre pair.
    \item La somme de trois nombres pairs consécutifs vaut 318. Quels sont ces nombres?
  \end{enumerate}
